\scanpage{11}
% archive sheet number: 10

\uline{Définition} Secteur polyédral : partie satisfaisant aux conditions équivalentes précédentes.

\uline{Prop} $E,E'$ vectoriels de dim finie en dualité. Alors par $A\unsure{\vdash|\mapsto}A^{\unsure{\circ|0}}$, $A'\vdash A'{}^\circ$ on trouve des correspondances $1$-$1$ inverses l'une de l'autre entre secteurs polyédraux de $E$ \unsure{passant|contenant} \deleted{\unsure{par}} \unsure{l'origine|d'origine}, et sect. \unsure{pol.} de $E'$ \unsure{contnt|contenant} l'origine.

\textcircled{\unsure{8|$\ast$}} Polarité et facettes

$E,E'$ Vectoriels de dim finie en dualité. Correspondance biunivoque entre sous-espaces affines $F$ de $E$ ne contenant pas l'origine, et sous espaces \inserted{affines} $F'$ de $E'$ ne contenant pas l'origine, on a

\margin{Notation $F^{\unsure{\vee|\circ|0}}$,}

\[
 \dim F+\dim F^\vee=n-1\qquad(n=\dim E=\dim E')
\]

Correspondance caractérisée par les propriétés a) \unsure{Renverser|Renverse} la relation d'inclusion

b) \unsure{À|$A'$} $x\in E-\{0\}$ correspond $\{x\}^\vee=H(x)^{\bullet}=(\{x\}^{\circ})^\bullet$
